% GENERATED FILE. Edit canonical lessons, then run npm run generate:books.
% canonical-source-hash: fnv1a64:3d8e290a3cbaf78e
% canonical-lessons: RU-C446-skver, RU-C446-vyveska, RU-C446-vitrina, RU-C446-vodostok, RU-C446-laryok

\chapter{Shop signs and windows}
\label{ch:ru-shop-signs-and-windows}
\hlchaptermodality{446}

\begin{chapteropening}
\textbf{By the end of this chapter:} \emph{I can describe shopfronts and small spaces: a little public garden, a shop sign, a shop window, a drainpipe and a stall.}

Complete the last lesson of chapter 446: I can describe shopfronts and small spaces: a little public garden, a shop sign, a shop window, a drainpipe and a stall.
\end{chapteropening}

\section[\texorpdfstring{\ru{сквер} (skver)}{сквер (skver)}]{\ru{сквер} (skver) — a small public garden, a square}

\label{lesson:RU-C446-skver}



\begin{quote}
Before the new one: say the Russian for a kiosk, then the Russian for a flower bed.
\end{quote}

\subsection*{You'll want to know: \ru{сквер}}
\textbf{\ru{сквер}} (\emph{skver}) — \textquotedblleft{}a small public garden, a square\textquotedblright{}.

\begin{quote}
\ru{Возле} \ru{театра} \ru{есть} \ru{маленький} \ru{сквер}.
\end{quote}

\begin{quote}
\emph{Vózle teátra yest málenkiy skver.}
\end{quote}

\begin{quote}
Near the theatre there is a small public garden.
\end{quote}

Masculine. A small green square with benches.

\begin{cousinweb}
\emph{\ru{сквер}} was borrowed from English \emph{square} (EtymDB); in Russian it names a small garden on a square.
\end{cousinweb}

\begin{grammarlens}[title={What you've built}]
One more word for telling, explaining and describing at B1.
\end{grammarlens}

\subsection*{Your turn}
\begin{itemize}
\raggedright
  \item \emph{Say it:} \emph{skver}
  \item \emph{Say it:} \emph{skver}, once more
  \item \emph{Say it:} a sentence of your own with \emph{skver}
  \item \emph{Recall:} say the Russian for a kiosk, then the Russian for a flower bed, then say \emph{skver} again
\end{itemize}

\begin{tcolorbox}[breakable,colback=teal!4,colframe=teal!35!black,title={Before you move on}]
What does \ru{сквер} mean? (\textquotedblleft{}A small public garden, a square\textquotedblright{}.) Say it once more.
\end{tcolorbox}
\section[\texorpdfstring{\ru{вывеска} (výveska)}{вывеска (výveska)}]{\ru{вывеска} (výveska) — a shop sign, a signboard}

\label{lesson:RU-C446-vyveska}



\begin{quote}
Before the new one: say the Russian for a flower bed, then the Russian for a small public garden.
\end{quote}

\subsection*{You'll want to know: \ru{вывеска}}
\textbf{\ru{вывеска}} (\emph{výveska}) — \textquotedblleft{}a shop sign, a signboard\textquotedblright{}.

\begin{quote}
\ru{Над} \ru{дверью} \ru{висит} \ru{новая} \ru{вывеска}.
\end{quote}

\begin{quote}
\emph{Nad dvéryu visít nóvaya výveska.}
\end{quote}

\begin{quote}
A new sign hangs over the door.
\end{quote}

Feminine; the stress is on \emph{\ru{вы}-} (\emph{vý-}).

\begin{cousinweb}
\emph{\ru{вывеска}} is built on \emph{\ru{вывесить}} 'to hang out', from \emph{\ru{вы}-} 'out' and \emph{\ru{весить}}, \emph{\ru{вешать}} 'to hang': a sign hung outside a shop. EtymDB does not list the word.
\end{cousinweb}

\begin{grammarlens}[title={What you've built}]
One more word for telling, explaining and describing at B1.
\end{grammarlens}

\subsection*{Your turn}
\begin{itemize}
\raggedright
  \item \emph{Say it:} \emph{výveska}
  \item \emph{Say it:} \emph{výveska}, once more
  \item \emph{Say it:} a sentence of your own with \emph{výveska}
  \item \emph{Recall:} say the Russian for a flower bed, then the Russian for a small public garden, then say \emph{výveska} again
\end{itemize}

\begin{tcolorbox}[breakable,colback=teal!4,colframe=teal!35!black,title={Before you move on}]
What does \ru{вывеска} mean? (\textquotedblleft{}A shop sign, a signboard\textquotedblright{}.) Say it once more.
\end{tcolorbox}
\section[\texorpdfstring{\ru{витрина} (vitrína)}{витрина (vitrína)}]{\ru{витрина} (vitrína) — a shop window}

\label{lesson:RU-C446-vitrina}



\begin{quote}
Before the new one: say the Russian for a small public garden, then the Russian for a shop sign.
\end{quote}

\subsection*{You'll want to know: \ru{витрина}}
\textbf{\ru{витрина}} (\emph{vitrína}) — \textquotedblleft{}a shop window\textquotedblright{}.

\begin{quote}
\ru{В} \ru{витрине} \ru{стоят} \ru{новые} \ru{туфли}.
\end{quote}

\begin{quote}
\emph{V vitríne stoyát nóvyye túfli.}
\end{quote}

\begin{quote}
There are new shoes in the shop window.
\end{quote}

Feminine. Also a glass display case in a museum.

\begin{cousinweb}
\emph{\ru{витрина}} is usually traced to French \emph{vitrine} 'a shop window, a display case', which EtymDB builds on \emph{vitre} 'a pane of glass'.
\end{cousinweb}

\begin{grammarlens}[title={What you've built}]
One more word for telling, explaining and describing at B1.
\end{grammarlens}

\subsection*{Your turn}
\begin{itemize}
\raggedright
  \item \emph{Say it:} \emph{vitrína}
  \item \emph{Say it:} \emph{vitrína}, once more
  \item \emph{Say it:} a sentence of your own with \emph{vitrína}
  \item \emph{Recall:} say the Russian for a small public garden, then the Russian for a shop sign, then say \emph{vitrína} again
\end{itemize}

\begin{tcolorbox}[breakable,colback=teal!4,colframe=teal!35!black,title={Before you move on}]
What does \ru{витрина} mean? (\textquotedblleft{}A shop window\textquotedblright{}.) Say it once more.
\end{tcolorbox}
\section[\texorpdfstring{\ru{водосток} (vodostók)}{водосток (vodostók)}]{\ru{водосток} (vodostók) — a drainpipe, a gutter}

\label{lesson:RU-C446-vodostok}



\begin{quote}
Before the new one: say the Russian for a shop sign, then the Russian for a shop window.
\end{quote}

\subsection*{You'll want to know: \ru{водосток}}
\textbf{\ru{водосток}} (\emph{vodostók}) — \textquotedblleft{}a drainpipe, a gutter\textquotedblright{}.

\begin{quote}
\ru{Во} \ru{время} \ru{дождя} \ru{из} \ru{водостока} \ru{бежит} \ru{вода}.
\end{quote}

\begin{quote}
\emph{Vo vrémya dozhdyá iz vodostóka bezhít vodá.}
\end{quote}

\begin{quote}
During the rain water runs out of the drainpipe.
\end{quote}

Masculine. \emph{\ru{Водосточная} \ru{труба}} — a drainpipe.

\begin{cousinweb}
\emph{\ru{водосток}} is built from \emph{\ru{вода}} 'water' and \emph{\ru{сток}} 'a draining off', on the root of \emph{\ru{течь}} 'to flow' (EtymDB lists \emph{\ru{водосток}}).
\end{cousinweb}

\begin{grammarlens}[title={What you've built}]
One more word for telling, explaining and describing at B1.
\end{grammarlens}

\subsection*{Your turn}
\begin{itemize}
\raggedright
  \item \emph{Say it:} \emph{vodostók}
  \item \emph{Say it:} \emph{vodostók}, once more
  \item \emph{Say it:} a sentence of your own with \emph{vodostók}
  \item \emph{Recall:} say the Russian for a shop sign, then the Russian for a shop window, then say \emph{vodostók} again
\end{itemize}

\begin{tcolorbox}[breakable,colback=teal!4,colframe=teal!35!black,title={Before you move on}]
What does \ru{водосток} mean? (\textquotedblleft{}A drainpipe, a gutter\textquotedblright{}.) Say it once more.
\end{tcolorbox}
\section[\texorpdfstring{\ru{ларёк} (laryók)}{ларёк (laryók)}]{\ru{ларёк} (laryók) — a stall, a booth}

\label{lesson:RU-C446-laryok}



\begin{quote}
Before the new one: say the Russian for a shop window, then the Russian for a drainpipe.
\end{quote}

\subsection*{You'll want to know: \ru{ларёк}}
\textbf{\ru{ларёк}} (\emph{laryók}) — \textquotedblleft{}a stall, a booth\textquotedblright{}.

\begin{quote}
\ru{В} \ru{ларьке} \ru{у} \ru{остановки} \ru{продают} \ru{пирожки}.
\end{quote}

\begin{quote}
\emph{V larké u ostanóvki prodayút pirozhkí.}
\end{quote}

\begin{quote}
The stall by the bus stop sells pies.
\end{quote}

Masculine; the \ru{ё} drops out: \emph{\ru{ларька}}, \emph{\ru{ларьки}}.

\begin{cousinweb}
\emph{\ru{ларёк}} is a diminutive of \emph{\ru{ларь}} 'a chest, a bin', which came through Old East Slavic \emph{larĭ} from Old Norse \emph{lárr} 'a chest, a drawer' (EtymDB).
\end{cousinweb}

\begin{grammarlens}[title={What you've built}]
One more word for telling, explaining and describing at B1.
\end{grammarlens}

\subsection*{Your turn}
\begin{itemize}
\raggedright
  \item \emph{Say it:} \emph{laryók}
  \item \emph{Say it:} \emph{laryók}, once more
  \item \emph{Say it:} a sentence of your own with \emph{laryók}
  \item \emph{Recall:} say the Russian for a shop window, then the Russian for a drainpipe, then say \emph{laryók} again
\end{itemize}

\begin{tcolorbox}[breakable,colback=teal!4,colframe=teal!35!black,title={Before you move on}]
What does \ru{ларёк} mean? (\textquotedblleft{}A stall, a booth\textquotedblright{}.) Say it once more.
\end{tcolorbox}
